\scanpage{10}

\begin{center}
U N I V E R S I T E\\
DES SCIENCES ET TECHNIQUES\\
DU LANGUEDOC\\
\textasciicircum\textasciicircum\textasciicircum\textasciicircum\textasciicircum\textasciicircum\textasciicircum\textasciicircum\textasciicircum\textasciicircum\textasciicircum\textasciicircum\textasciicircum
\end{center}

Montpellier, le 18 octobre 1983

Service de la Scolarité

Diffusion : Liste A\\
AM/gs \hfill N° 4570

\begin{center}
AVIS DE SOUTENANCE D'UNE THESE DE DOCTEUR-INGENIEUR

\uuline{``ECOLOGIE GENERALE ET APPLIQUEE option ECOLOGIE TERRESTRE''}
\end{center}

\uline{NOM et prénom du candidat} : MOUFADDAL Mohamed

\uline{Titre de la thèse} : ``CONTRIBUTION A L'ETUDE DES POSSIBILITES D'EXTENSION DU CEDRE DE L'ATLAS EN DEHORS DE SON AIRE ACTUELLE''

\uline{Date, heure et lieu de soutenance prévus} : Le 28 octobre 1983 à 14 heures à l'Institut de Botanique

\uline{Nom du directeur de thèse} : M. GODRON Michel

\uline{Désignation du laboratoire où la thèse est préparée} : Centre d'Etudes Phytosociologiques et écologiques Louis Emberger (CNRS)

\uline{Aanlyse de l'ouvrage} : L'objet de cette étude étant de cerner les possibilités d'extension du Cèdre (\uline{Cedrus atlantica} Manetti) en dehors de son aire actuelle, il a fallu, après une analyse de l'histoire de répartition du Cèdre, préciser les facteurs climatiques responsables de la distribution actuelle.

Les régions où l'introduction du Cèdre est possible, sont déterminées à partir de l'étude climatique de stations situées à l'intérieur et à l'extérieur de l'aire du Cèdre, et par une étude écologique et dendrométrique comparative de quelques plantations.

\uline{Durée approximative des recherches ayant abouti à la thèse} : 5 ans

\uline{Composition du jury proposé} : \begin{tabular}[t]{ll}
MM. & GODRON M.\\
& BEDEL J.\\
& SCHOENENBERGER A.\\
& ROMANE F.\\
& DAGET Ph.
\end{tabular}

\begin{flushright}
P/ le Secrétaire Général,\\
l'Attachée principale d'administration\\
scolaire et universitaire,

\unsure{A. Gentil}

\uline{A. GENTIL}
\end{flushright}

\margin{9}
